\begin{abstract}
\porozhdeno{\Claim{A02}{Выбор референсной базы и связанной с ней предобученной модели представляет собой потенциальный источник вариабельности таксономических назначений при анализе ампликонных данных. Сравнение таких моделей при фиксированных условиях обработки помогает оценить устойчивость таксономических профилей и сопоставимость результатов исследований.}}
\porozhdeno{\Claim{A01}{Сравнены готовые модели, связанные с базами SILVA, GTDB и Greengenes2, при использовании единого алгоритма таксономического назначения и одинаковой обработки ампликонных данных почв Москвы. Для \ProjectSamplesSoil{} почвенных образцов оценены покрытие таксономическими назначениями, буквальное совпадение таксономических линий и расстояния между таксономическими профилями. Медиана расстояния между профилями SILVA и GTDB на уровне семейств составила \CmpSilvaGtdbFamilyDistanceMedian{} при межобразцовом диапазоне \ArticleSilvaGtdbFamilyDistanceMin{}--\ArticleSilvaGtdbFamilyDistanceMax{}; при построении профилей учтена неназначенная доля. Медианы покрытия ридов назначениями на уровне семейства составили \CmpSilvaFamilyReadMedianPct{}\% и \CmpGtdbFamilyReadMedianPct{}\% для SILVA и GTDB соответственно. Отбор бактерий и архей отдельно по назначениям каждой модели изменяет знаменатели рассчитываемых показателей и может приводить к отсутствию отобранных последовательностей либо к отсутствию назначений на уровне семейства при непустом наборе. Численные значения проверены независимым скриптом по сохранённым результатам. Выявленные различия характеризуют чувствительность таксономических назначений к выбору готовой модели. В отсутствие эталона они не позволяют оценить точность классификации и сами по себе не служат основанием для экологической интерпретации состава почвенных сообществ.}}
\end{abstract}
